% concatenated from arxiv2_r218.tex
\documentclass{amsart}

%\usepackage[a4paper]{geometry}
%\usepackage[letterpaper]{geometry}
%\usepackage[textwidth=7in]{geometry}
\usepackage[margin=1in]{geometry}


\newtheorem{theorem}{Theorem}[section]
\newtheorem{lemma}[theorem]{Lemma}

\theoremstyle{definition}
\newtheorem{definition}[theorem]{Definition}
\newtheorem{example}[theorem]{Example}
\newtheorem{xca}[theorem]{Exercise}

\theoremstyle{remark}
\newtheorem{remark}[theorem]{Remark}
\numberwithin{equation}{section}

\usepackage{booktabs}
\usepackage{siunitx}
\usepackage{graphicx}
\usepackage{hyperref}
\usepackage{pdflscape}
\usepackage{xcolor}


\newtheorem{corollary}{Corollary}[theorem] % Use theorem counter as `parent`
\newtheorem{proposition}[theorem]{Proposition}

% \documentclass{article}
\usepackage{listings}
\usepackage{xcolor}

\lstset{
	language=R,
	basicstyle=\ttfamily\small,
	backgroundcolor=\color{gray!10},
	frame=single,
	keywordstyle=\color{blue},
	commentstyle=\color{green!40!black},
	stringstyle=\color{orange},
	showstringspaces=false
}


\begin{document}

\title{On the Bounds of Stopping and Cycle Numbers in the Collatz Process}

%    Only \author and \address are required; other information is
%    optional.  Remove any unused author tags.

%    author one information
% \author[short version for running head]{name for top of paper}
\author{Daohang Sha}
\address{}
\curraddr{Philadelphia, PA 19104}
\email{dhsha@hotmail.com}
\thanks{}

%    author two information
%\author{}
%\address{}
%\curraddr{}
%\email{}
%\thanks{}

%    \subjclass is required.
\subjclass[2020]{Primary 11Y16, 00A05; Secondary 33D99, 68Q25}
\keywords{Collatz process, 3n+1 problem, Stopping time, Cycle number}
\date{Jan 24, 2026}

%\dedicatory{}

%    Abstract is required.

% short version
\begin{abstract}
	We analyze the stopping-time and cycle structure of the normalized Collatz
	iteration. Using a recursive description of admissible binary sequences, we show
	that every integer $m \equiv 3 \pmod{4}$ arises uniquely and derive new bounds for
	the associated stopping and cycle numbers. These bounds imply that
	$F_{q}(m)/m \to 1$ as the sequence length increases, while equality is impossible
	for any finite sequence. Consequently, no finite nontrivial cycle is compatible
	with the iteration, and the trivial cycle at $1$ is the only admissible periodic
	orbit.
\end{abstract}


\maketitle


\sloppy

\section{Introduction}
The $3n+1$ problem, also known as the Collatz conjecture, has been widely studied and remains one of the most intriguing unsolved problems in mathematics \cite{La11}, \cite{La12}. While several formulations of the Collatz function exist, we consider the following normalized form:
$$f(n)=\begin{cases}
	F_0(n)=\frac{n}{2}, & n = even \\
	F_1(n)=\frac{3n+1}{2}, & n = odd
\end{cases}$$

We denote the Collatz process as $F_q(n)=F_{101...0}(n)=F_0(...F_1(F_0(F_1(n))))$, where the binary sequence $q=101...0$, the total steps or the length of sequence is $s$ and the total number of odd steps is $r$. If $F_q(n)<n$, the minimum step $s$ is called the stopping time. If $F_q(m_c)=m_c$, $m_c$ is called a cycle number. The stopping time ($s$) and the cycle numbers ($m_c$) have been widely studied in \cite{Si10} - \cite{He23}. In \cite{Si10}, Simons and de Weger referred to a periodic sequence with $m$ local minima as an $m-cycle$. In the following discussion, we use the term "cycle" to denote any sequence that contains at least one such cycle. We also restrict our analysis to odd starting values $n \equiv 3 \pmod{4}$ satisfying $s \geq 4$, as those with $s \leq 2$ are readily confirmed by direct computation.

In this paper we obtain a general lower bound for the stopping number and compute sharp upper bounds for both stopping and cycle numbers across all admissible sequences. Taken together, these bounds show that any nontrivial cycle would require an exact identity between powers of 2 and 3 in the expression for $F_q(n)$, an equality unattainable for any finite sequence. Thus no finite nontrivial cycle can occur in the normalized Collatz iteration.

\section{Collatz Sequence of Stopping}

Let $q$ be a binary sequence of length $s$, containing exactly $r$ entries equal to 1. We define the function $1(i)$ to denote the position of the $i^{\text{th}}$ occurrence of the digit 1 within the sequence $q$. For example, consider $n=423$, $F_q(n)=302$, the sequence $q = 1110110100$, which contains six ones. The corresponding positions of these ones are given by $1(1) = 1$, $1(2) = 2$, $1(3) = 3$, $1(4) = 5$, $1(5) = 6$, and $1(6) = 8$. Using this notation, the transformation $F_q(n)$ associated with the sequence $q$ in the Collatz process is given by
\begin{equation} \label{formula}
	F_q (n)=\dfrac{3^r n+\sum_{i=1}^r3^{r-i} 2^{1(i)-1}}{2^s},
\end{equation}

Here, $s$ is the total number of steps, $r$ is the number of applications of the function $F_1$, and $1(i)$ specifies the position of the $i^{\text{th}}$ occurrence of 1 in the sequence $q$.


\begin{lemma}\label{lem:length}
	Let \( r \in \mathbb{N} \) and define
	\[
	s = \left\lfloor 1 + r \cdot \log_2 3 \right\rfloor.
	\]
	Then there exist a non-negative integer \( i \in \mathbb{Z}_{\geq 0} \) and a positive integer \( k \in \mathbb{N} \) such that
	\[
	3^{r} \cdot (2i+1) - 1 = k \cdot 2^{s - r}.
	\]
\end{lemma}

\begin{proof}
	By the definition of \( s \), we have
	\[
	s = \left\lfloor 1 + r \cdot \log_2 3 \right\rfloor \quad \Rightarrow \quad s - 1 < r \log_2 3 < s.
	\]
	Exponentiating both sides with base 2 gives
	\[
	2^{s - 1} < 3^r < 2^s,
	\]
	which implies
	\[
	\frac{3^r}{2^{s - r}} \in \left(2^{r - 1}, 2^r\right).
	\]
	
	Now consider the congruence
	\[
	3^r \cdot t \equiv 1 \pmod{2^{s - r}}.
	\]
	Since \( 3^r \) is odd and \( \gcd(3^r, 2^{s - r}) = 1 \), the inverse of \( 3^r \) modulo \( 2^{s - r} \) exists. Therefore, there exists an odd integer \( t \equiv 2i + 1 \), for some \( i \in \mathbb{Z}_{\geq 0} \), such that
	\[
	3^r \cdot (2i + 1) \equiv 1 \pmod{2^{s - r}}.
	\]
	Rewriting this congruence yields
	\[
	3^r \cdot (2i + 1) - 1 = k \cdot 2^{s - r}
	\]
	for some \( k \in \mathbb{N} \), as required.
\end{proof}


\begin{lemma}\label{lem:min}
	Let \( r\geq2, i \in \mathbb{N} \), and for any starting number 
	\[
	m = 2^r \cdot (2i+1) - 1 \equiv 3\pmod{4}.
	\]
	Then the Collatz sequence of \( m \) begins with
	\[
	\underbrace{11\ldots1}_{r \text{ times}}\,\underbrace{00\ldots0}_{s - r \text{ times}}.
	\]
\end{lemma}

\begin{proof}
	Let us begin with the expression
	\[
	m = 2^r \cdot (2i+1) - 1.
	\]
	
	We claim that applying \( r \) iterations of the Collatz rule for odd integers, i.e., \( m \mapsto (3m + 1)/2 \), leads to
	\[
	F_{11...1}(m) = 3^{r} \cdot (2i+1) - 1.
	\]
	
	This is shown by induction:
	
	- Base case (step 1):
	\[
	\quad F_1(m) = \frac{3(2^r \cdot (2i+1) - 1) + 1}{2} = 3 \cdot 2^{r-1} (2i+1) - 1 = 2^{r-1} \cdot 3^{1} \cdot (2i+1) - 1.
	\]
	
	- Inductive step:
	Suppose that after \( r-1 \) steps,
	\[
	F_{11...1}(m) = 2 \cdot 3^{r - 1} \cdot (2i+1) - 1.
	\]
	Then, step $r$:
	\[
	F_{11...11}(m) = \frac{3[2 \cdot 3^{r - 1} \cdot (2i+1) - 1] + 1}{2} = \frac{2 \cdot 3^{r} \cdot (2i+1) - 2}{2} = 3^{r} \cdot (2i+1) - 1.
	\]
	
	So, after \( r \) such "odd steps" (each denoted by `1`), we reach \( 3^{r} (2i+1) - 1 \), which is even (since \( 3^{r} \) is odd, so the result is even). From this point, the Collatz process proceeds with divisions by 2 only, i.e., steps of the form \( m \mapsto m/2 \), which we denote by '0'.
	
	Hence, the step sequence begins with exactly \( r \) increasing steps ('1's), followed by decreasing steps ('0's) until the next odd number occurs.
	
	Therefore, together with the Lemma~\ref{lem:length}, the sequence begins with:
	\[
	\underbrace{11\ldots1}_{r \text{ times}}\,\underbrace{00\ldots0}_{s - r \text{ times}},
	\]
	completing the proof.
\end{proof}

\tiny
%\scriptsize
%\footnotesize
%\small
\begin{table}[ht]
	\caption{Specific Sequences of $q$ with $1(i)=i$ for $m \equiv 3\pmod{4}$ in Collatz process}
	\label{tab:seq}
\begin{tabular}{p{0.005\textwidth}p{0.005\textwidth}p{0.4\textwidth}p{0.11\textwidth}rp{0.125\textwidth}p{0.125\textwidth}}

		\toprule

		$r$ &  {$s$} & {Sequence($q$)} & {$2^r\cdot(2i+1)-1$} & {$m\pmod{2^s}$} & {$k=F_q(m)$} \\

		\midrule
		2	&	4	&	1100 & {$2^2 \cdot 1-1$} & 3 & 2\\
		3	&	5	&	11100 & {$2^3 \cdot 3-1$} & 23 & 20\\
		4	&	7	&	1111000 & {$2^4 \cdot 1-1$} & 15 & 10\\
		5	&	8	&	11111000 & {$2^5 \cdot 3-1$} & 95 & 91\\
		6	&	10	&	1111110000 & {$2^6 \cdot 9-1$} & 575 & 410\\
		7	&	12	&	111111100000 & {$2^7 \cdot 3-1$} & 383 & 205\\
		8	&	13	&	1111111100000 & {$2^8 \cdot 1-1$} & 255 & 205\\
		9	&	15	&	111111111000000	& {$2^9 \cdot 11-1$} & 5,631 & 3,383\\
		10	&	16	&	1111111111000000 & {$2^{10} \cdot 25-1$} & 25,599 & 23,066\\
		11	&	18	&	111111111110000000 & {$2^{11} \cdot 51-1$} & 104,447 & 70,582\\
		12	&	20	&	11111111111100000000 & {$2^{12} \cdot 17-1$} & 69,631 & 35,291\\
		13	&	21	&	111111111111100000000 & {$2^{13} \cdot 91-1$} & 745,471 & 566,732\\
		14	&	23	&	11111111111111000000000 & {$2^{14} \cdot 201-1$} & 3,293,183 & 1,877,689\\
		15	&	24	&	111111111111111000000000 & {$2^{15} \cdot 67-1$} & 2,195,455 & 1,877,689\\
		16	&	26	&	11111111111111110000000000 & {$2^{16} \cdot 193-1$} & 12,648,447 & 8,113,298\\
		17	&	27	&	111111111111111110000000000 & {$2^{17} \cdot 747-1$} & 97,910,783 & 94,206,740\\
		18	&	29	&	11111111111111111100000000000 & {$2^{18} \cdot 249-1$} & 65,273,855 & 47,103,370\\
		19	&	31	&	1111111111111111111000000000000	& {$2^{19} \cdot 83-1$}  & 43,515,903 & 23,551,685\\
		20	&	32	&	11111111111111111111000000000000 & {$2^{20} \cdot 1393-1$} & 1,460,666,367 & 1,185,813,152\\
		21	&	34	&	1111111111111111111110000000000000 & {$2^{21} \cdot 3195-1$} & 6,700,400,639 & 4,079,690,977\\
		22	&	35	&	11111111111111111111110000000000000 & {$2^{22} \cdot 1065-1$} & 4,466,933,759 & 4,079,690,977\\
		23	&	37	&	1111111111111111111111100000000000000 & {$2^{23} \cdot 8547-1$} & 71,697,432,575 & 49,111,434,902\\
		24	&	39	&	111111111111111111111111000000000000000 & {$2^{24} \cdot 2849-1$} & 47,798,288,383 & 24,555,717,451\\
		25	&	40	&	1111111111111111111111111000000000000000 & {$2^{25} \cdot 22795-1$} & 764,873,277,439 & 589,414,790,413\\
		26	&	42	&	111111111111111111111111110000000000000000 & {$2^{26} \cdot 18521-1$} & 1,242,923,270,143 & 718,351,699,928\\
		27	&	43	&	1111111111111111111111111110000000000000000 & {$2^{27} \cdot 28019-1$} & 3,760,646,520,831 & 3,260,217,528,257\\
		28	&	45	&	111111111111111111111111111100000000000000000 & {$2^{28} \cdot 31185-1$} & 8,371,159,695,359 & 5,442,907,506,622\\
		29	&	46	&	1111111111111111111111111111100000000000000000 & {$2^{29} \cdot 10395-1$} & 5,580,773,130,239 & 5,442,907,506,622\\
		30	&	48	&	111111111111111111111111111111000000000000000000 & {$2^{30} \cdot 3465-1$} & 3,720,515,420,159 & 2,721,453,753,311\\
		31	&	50	&	11111111111111111111111111111110000000000000000000 & {$2^{31} \cdot 263299-1$} & 565,430,297,034,751 & 310,197,425,018,629\\	
		
		\vdots	&	\vdots	&	\vdots & \vdots &\vdots &	\vdots 	\\	
		
		\bottomrule

	\end{tabular}
%	}
\end{table}
\normalsize

Table~\ref{tab:seq} includes several specific sequences for values of $r$ ranging from 2 to 31. This list can be extended indefinitely, as there is no upper bound on $r$. Therefore, there is no upper bound on $s$, too. For example, when $r=1000$, then $s=1585$,\\

$m=113181711197902800523896881066218160102709667045393956950389619432948098013957729384098\\
0894314243553691246663949249137551627877867250646145606647381546299358516204398946163838768783\\
5954105697727172101185270096853888059974345842688019492876298064702543347682952175404153595977\\
9733727568209962344081540465463827583787623071447065173716333421030611964595317030946266729588\\
1872408781238846687069393680025440093827021268188760756952836484710583975703867188025783792797\\
611927080534015$ (478 digits)\\

$F_q(m)=11027773766223939255599638243685383983123058040229118529895194635849895633361467411\\
1382402477750726619658150157035536038452913603129669634101848918205329307515073658919130294253\\
2781506097769243834069629788450644321665401676961180452393456153975938038355963675587842309211\\
9591171602528480255231243448888268364113078890677027695508662577537129844756989406360976065975\\
1963882667459405174174888685722485109676483971192668155890160044032425213449698721946740600735\\
2136396783848666645$ (478 digits)\\

$2i+1=105628373340941055460453010826461801731312915456661368786503022090718002344954512657\\4
11550899143065854323968254783269393108263739680866504502238016353334720601889819555845564641$ (177 digits)\\

\begin{theorem}\label{thm:min}
	The summation $ \sum_{i=1}^r 3^{r-i} \cdot 2^{1(i)-1} \geq 3^{r} - 2^{r},$ and equality holds when $1(i) = i$ for all $i$.
\end{theorem}

\begin{proof}
	To find the minimum value of
	$	\sum_{i=1}^r 3^{r-i} \cdot 2^{1(i)-1} $
	where $1(i)$ is the position of the $i^\text{th}$ 1 in a binary sequence $q$ of length $s$, we apply the result from the Rearrangement Inequality.	The minimum occurs when the positions $1(i)$ are as small as possible, i.e., $	1(i) = i \quad \text{for } i = 1, 2, \dots, r $
	That is, the 1s occur at the first $r$ positions in the sequence. Therefore, the minimum value is:
	$$	\sum_{i=1}^r 3^{r-i} \cdot 2^{i-1}	$$
    We can compute this closed-form as a geometric convolution:
	$$
	\sum_{i=1}^r 3^{r-i} \cdot 2^{i-1} = \sum_{j=0}^{r-1} 3^{j} \cdot 2^{r-1 - j} = 2^{r-1} \sum_{j=0}^{r-1} \left( \frac{3}{2} \right)^j
	$$
	This is a geometric series:	
	$$
	= 2^{r-1} \cdot \frac{ \left( \frac{3}{2} \right)^r - 1 }{ \frac{3}{2} - 1 } = 2^{r-1} \cdot \frac{ \left( \frac{3}{2} \right)^r - 1 }{ \frac{1}{2} } = 2^r \left[ \left( \frac{3}{2} \right)^r - 1 \right] = 3^{r} - 2^{r}
	$$
	Final Answer:
	$$
	\boxed {\sum_{i=1}^r 3^{r-i} \cdot 2^{1(i)-1} \geq 3^{r} - 2^{r} } 
	$$
	and equality holds when $1(i) = i$ for all $i$.
	
\end{proof}

\begin{theorem}\label{thm:max}
	The summation 
	\[
	\sum_{i=1}^r 3^{\,r-i} \cdot 2^{\,1(i)-1}
	\]
	attains its maximum value when 
	\[
	1(i) = \left\lfloor 1 + (i-1)\log_2 3 \right\rfloor, \quad i=1,\dots,r.
	\]
\end{theorem}

\begin{proof}
	Consider the modified rooted sequence tree introduced by Winkler \cite{Win17}, illustrated in  \autoref{fig:seqtree}. It is evident that the further to the right the $1$’s appear in a sequence, the larger the resulting sum becomes comparing to the minimum summation when all $1$’s appear in the left. Consequently, the sequences along the top line of the sequence tree yield the maximum values.
		
	\begin{figure}[ht]
		\centering
		\includegraphics[width=0.99\linewidth]{sequence_tree}
		\caption{Rooted sequence tree}
		\label{fig:seqtree}
	\end{figure}

Here is how these sequences in the modified sequence tree are constructed. All sequences are begin with $r=2, s=4$, so that the initial sequence is $q=1100$. This is ao called rooted tree. Each sub-sequence is obtained by inserting either “10” or “1” immediately before each zero at the terminal zero-block in the previous sequence, according to the difference in their lengths. Specifically, define
\[
d(r) = s(r) - s(r-1) = 
\left\lfloor 1 + r \log_2 3 \right\rfloor - 
\left\lfloor 1 + (r-1)\log_2 3 \right\rfloor.
\]

This rule generates all admissible sub-sequences. And the total number of the admissible sequences in each $r$ equal to the summary of all number of zeros at the terminal zero-block in the previous sequence. Consequently, the positions of the (1)’s in the top-line sequence are given exactly by
\[
1(i) = \left\lfloor 1 + (i-1)\log_2 3 \right\rfloor, 
\quad i=1,\dots,r.
\]

This establishes the claim.

\end{proof}

Table~\ref{tab:seq_max} shows the Collatz process of $r$ from 2 to 31, but with sequences with the maximum values of the summation. Again, the value of $r$ can be any big integers.

\tiny
%\scriptsize
%\footnotesize
%\small
\begin{table}[ht]
	\caption{Specific Sequences of $q$ with max value of the sum for $m \equiv 3\pmod{4}$ in Collatz process}
	\label{tab:seq_max}
		\begin{tabular}{p{0.005\textwidth}p{0.005\textwidth}p{0.4\textwidth}rp{0.13\textwidth}p{0.13\textwidth}}

				
			\toprule
			%		&   & {$N$} &  &  & {$4i+3\equiv 3\pmod{4}  $} & \\
			$r$ &  {$s$} & {Sequence($q$)} & {$m\pmod{2^s}$} &  {$F_q(m)<m$} \\
			%		&  {$\equiv0(mod3)$} & {$\equiv2(mod3)$} & {$\equiv1(mod3)$} & {$\equiv3(mod12)$} & {$\equiv7(mod12)$} & {$\equiv11(mod12)$}\\
			\midrule
			2 &	4	& 1100     						& 3			&	2\\
			3 &	5	& 11010							& 11		&	10\\
			4 &	7	& 1101100						& 59		&	38\\
			5 &	8	& 11011010      				& 123		&	118\\
			6 &	10	& 1101101100	  				& 507		&	362\\
			7 &	12	& 110110110100  				& 1,019 	&	545\\
			8 &	13	& 1101101101010		 			& 3,067		&	2,458\\
			9 &	15	& 110110110101100 				& 7,163  	&	4,304\\
			10 & 16	& 1101101101011010 				& 56,315	&	50,743\\
			11 & 18	& 110110110101101100 	    	& 23,547	&	15,914\\
			12 & 20	& 11011011010110110100    		& 416,763  	&	211,226\\
			13 & 21	& 110110110101101101010   		& 941,051 	&	715,420\\
			14 & 23	& 11011011010110110101100  		& 8,281,083	&	4,721,663\\
			15 & 24	& 110110110101101101011010		& 4,086,779 &	3,495,268\\
			16 & 26	& 11011011010110110101101100	& 62,807,035 &	40,287,332\\
			17 & 27	& 110110110101101101011011010	& 96,361,467 &  92,716,039\\
			18 & 29	& 11011011010110110101101101100	& 431,905,787 & 311,674,835\\
			19 & 31	& 1101101101011011010110110110100 & 163,470,331 & 88,473,443\\
			20 & 32	& 11011011010110110101101101101010 & 1,237,212,155 & 1,004,406,265\\
			21 & 34	& 1101101101011011010110110110101100 & 7,679,663,099 & 4,675,937,150\\
			22 & 35	& 11011011010110110101101101101011010 &	33,449,466,875 & 30,549,700,432\\
			23 & 37	& 1101101101011011010110110110101101100	& 119,348,812,795 & 81,751,762,091\\
			24 & 39	& 110110110101101101011011011010110110100 &	462,946,196,475 & 237,832,281,869\\
			25 & 40	& 1101101101011011010110110110101101101010 & 188,068,289,531 & 144,926,270,443\\
			26 & 42	& 110110110101101101011011011010110110101100 & 2,936,847,358,971 & 1,697,360,845,538\\
			27 & 43	& 1101101101011011010110110110101101101011010 &	5,135,870,614,523 & 4,452,440,639,554\\
			28 & 45	& 110110110101101101011011011010110110101101100 & 18,330,010,147,835 & 11,918,127,650,276\\
			29 & 46	& 1101101101011011010110110110101101101011011010 & 35,922,196,192,251 & 35,034,785,816,635\\
			30 & 48	& 110110110101101101011011011010110110101101101100 & 211,844,056,636,411 & 154,958,046,921,632\\
			31 & 50	& 11011011010110110101101101101011011010110110110100 & 915,531,498,413,051 & 502,264,407,868,691\\
																
			\vdots	&	\vdots	&	\vdots & \vdots &\vdots 	\\	
			\bottomrule
			%\end{tabular}
		\end{tabular}
%	}
	\end{table}
\normalsize


\begin{lemma}\label{lem:pow2}
	For any integer $s \ge 1$,
	\[
	\sum_{i=0}^{s-1} 2^{\,i} = 2^{\,s} - 1.
	\]
\end{lemma}

\begin{proof}
	Consider the finite geometric series with ratio $2$:
	\[
	\sum_{i=0}^{s-1} 2^{\,i} = \frac{2^{\,s} - 1}{2 - 1} = 2^{\,s} - 1.
	\]
\end{proof}


\begin{lemma}\label{lem:3mod4}

Let $m \equiv 3 \pmod{4}$. Then $m$ admits a unique representation of the form
$$ m = 3 + \sum_{n_i\ge2} 2^{n_i}, \qquad n_i \ge 2, $$
where the exponents $n_i$ are distinct. Conversely, every finite sum of the form
$$ 3 + \sum_{n_i\ge2} 2^{n_i}, \qquad n_i \ge 2,$$
yields an integer $m \equiv 3 \pmod{4}$.

\end{lemma}

\begin{proof}

Since $2^{n} \equiv 0 \pmod{4}$ for all $n \ge 2$, any integer of the form $3 + \sum_{n_i} 2^{n_i}$ is congruent to $3\pmod{4}$. Conversely, if $m \equiv 3 \pmod{4}$, then $m-3$ is divisible by 4 and hence has a binary expansion involving only powers $2^{n}$ with $n \ge 2$. Uniqueness follows from the uniqueness of the binary expansion.

\end{proof}


\begin{example}	
For instance,
%\tiny
%\scriptsize
%\footnotesize
\small
\[
\begin{array}{c|l}
	4i-1 \equiv 3\pmod{4} & \sum_{n_i \ge 2} 2^{n_i} + 3 \\ \hline
	3  & 3 \\
	7  & 2^{2} + 3 \\
	11 & 2^{3} + 3 \\
	15 & 2^{3} + 2^{2} + 3 \\
	19 & 2^{4} + 3 \\
	23 & 2^{4} + 2^{2} + 3 \\
	27 & 2^{4} + 2^{3} + 3 \\
	31 & 2^{4} + 2^{3} + 2^{2} + 3 \\
	35 & 2^{5} + 3 \\
	39 & 2^{5} + 2^{2} + 3 \\
	43 & 2^{5} + 2^{3} + 3 \\
	47 & 2^{5} + 2^{3} + 2^{2} + 3 \\
	51 & 2^{5} + 2^{4} + 3 \\
	55 & 2^{5} + 2^{4} + 2^{2} + 3 \\
	59 & 2^{5} + 2^{4} + 2^{3} + 3 \\
	63 & 2^{5} + 2^{4} + 2^{3} + 2^{2} + 3 \\
	67 & 2^{6} + 3 \\
	71 & 2^{6} + 2^{2} + 3 \\
	75 & 2^{6} + 2^{3} + 3 \\
	79 & 2^{6} + 2^{3} + 2^{2} + 3 \\
	83 & 2^{6} + 2^{4} + 3 \\
	87 & 2^{6} + 2^{4} + 2^{2} + 3 \\
	91 & 2^{6} + 2^{4} + 2^{3} + 3 \\
	95 & 2^{6} + 2^{4} + 2^{3} + 2^{2} + 3\\
	\vdots & \vdots
\end{array}
\]
\normalsize
\end{example}


\begin{lemma}[Uniqueness of sequences and starting values]\label{lem:seqtree}
	The binary sequences $q$ displayed in Figure~\ref{fig:seqtree} are pairwise distinct. 
	Moreover, the construction defines an injective mapping
	\[
	q \;\longmapsto\; m,
	\]
	where each sequence $q$ corresponds to a unique integer $m \equiv 3 \pmod{4}$
	satisfying $F_q(m) < m$.
\end{lemma}

\begin{proof}
	The sequences in Figure~\ref{fig:seqtree} are generated recursively using a
	deterministic insertion rule: at each stage $r$, a new sequence is obtained from a
	unique parent sequence by inserting either the substring $1$ or $10$ at a prescribed
	position immediately preceding the terminal block of zeros. Since both the parent
	sequence and the insertion position are uniquely determined, distinct construction
	paths necessarily produce distinct binary sequences. Hence all sequences $q$ in
	Figure~\ref{fig:seqtree} are pairwise distinct.
	
	For a fixed sequence $q$ of length $s$, equation~\eqref{formula} uniquely determines
	the residue class of $m \pmod{2^s}$. By Lemma~\ref{lem:3mod4}, every integer
	$m \equiv 3 \pmod{4}$ admits a unique binary expansion of the form
	\[
	m = 3 + \sum_{n_i \ge 2} 2^{n_i}.
	\]
	The exponents $n_i$ correspond uniquely to the positions of the $1$'s in the
	sequence $q$. Consequently, distinct sequences $q$ give rise to distinct integers
	$m$. Finally, direct comparison of the iterates of $m$ with its binary expansion shows
	that the constructed integer satisfies $F_q(m) < m$. This establishes injectivity
	of the mapping $q \mapsto m$.
\end{proof}

\begin{example}	
For instance, let $i$ is any integer and $4i-1 \equiv 3\pmod{4}$, then the first two iterates are
\[
F_{1}(4i-1) = 3\cdot 2i - 1, 
\qquad 
F_{11}(4i-1) = 3^2 i - 1.
\]

As an illustrative example, take 
\[
r = 11, \quad s = 18, \quad q = 111111111010100000.
\]
Then the iterates proceed as
\[
\begin{aligned}
	F_{111}(4i-1) &= \frac{3(3^2(2j)-1)+1}{2} = 3^3 j - 1,\\
	F_{1111}(4i-1) &= \frac{3(3^3(2k)-1)+1}{2} = 3^4 k - 1,\\
	F_{11111}(4i-1) &= \frac{3(3^4(2\ell)-1)+1}{2} = 3^5 \ell - 1,\\
	F_{111111}(4i-1) &= \frac{3(3^5(2m)-1)+1}{2} = 3^6 m - 1,\\
	F_{1111111}(4i-1) &= \frac{3(3^6(2n)-1)+1}{2} = 3^7 n - 1,\\
	F_{11111111}(4i-1) &= \frac{3(3^7(2o)-1)+1}{2} = 3^8 o - 1,\\
	F_{111111111}(4i-1) &= \frac{3(3^8(2p)-1)+1}{2} = 3^9 p - 1,\\
	F_{1111111110}(4i-1) &= \frac{3^9(2q+1)-1}{2} = 3^9 q + 9841,\\
	F_{11111111101}(4i-1) &= \frac{3(3^9(2r)+9841)+1}{2} = 3^{10} r + 14762,\\
	F_{111111111010}(4i-1) &= \frac{3^{10}(2s)+14762}{2} = 3^{10} s + 7381,\\
	F_{1111111110101}(4i-1) &= \frac{3(3^{10}(2t)+7381)+1}{2} = 3^{11} t + 11072,\\
	F_{11111111101010}(4i-1) &= \frac{3^{11}(2u)+11072}{2} = 3^{11} u + 5536,\\
\end{aligned}	
\]
\[
\begin{aligned}
	F_{111111111010100}(4i-1) &= \frac{3^{11}(2v)+5536}{2} = 3^{11} v + 2768,\\
	F_{1111111110101000}(4i-1) &= \frac{3^{11}(2w)+2768}{2} = 3^{11} w + 1384,\\
	F_{11111111101010000}(4i-1) &= \frac{3^{11}(2x)+1384}{2} = 3^{11} x + 692,\\
	F_{111111111010100000}(4i-1) &= \frac{3^{11}(2y)+692}{2} = 3^{11} y + 346 = 3^r y + F_q(m).
\end{aligned}
\]

The corresponding expansion of $4i-1$ in powers of two is
\[
\begin{aligned}
	4i-1 &= 2^2 (2j) - 1 = 2^2 2^2 k - 1 = 2^2 2^3 \ell - 1 = 2^2 2^4 m - 1 = 2^2 2^5 n - 1 = 2^2 2^6 o - 1\\
	&= 2^2 2^7 p - 1 = 2^2 2^7 (2q+1) - 1 = 2^2 2^8 q + 2^2 2^7 - 1 = \hdots = 2^{18} y + 511\\
	&= 2^s y + \sum_{n_i \ge 2} 2^{n_i} + 3 = 2^s y + m.\\
\end{aligned}
\]

\end{example}


\begin{theorem}[Exhaustivity for $m \equiv 3 \pmod{4}$]\label{thm:exhaustive}
	The construction in Figure~\ref{fig:seqtree} generates each integer
	$m \equiv 3 \pmod{4}$ exactly once. Equivalently, the correspondence
	\[
	q \;\longleftrightarrow\; m
	\]
	defined by Figure~\ref{fig:seqtree} is a bijection between admissible sequences $q$
	and integers $m \equiv 3 \pmod{4}$.
\end{theorem}

\begin{proof}
	We prove the statement by showing that the constructions
	$q \mapsto m$ and $m \mapsto q$ are well-defined and inverse to each other.
	
	\medskip
	\noindent\textbf{Forward direction ($q \mapsto m$).}
	Starting from the root $q = 11$, we have
	\[
	F_{11}(4i - 1) = 3^{2} i - 1.
	\]
	Any admissible sequence $q = 11\cdots$ corresponds to a path in the rooted
	sequence tree shown in Figure~\ref{fig:seqtree}. Along this path, the Collatz
	process, $F_q(4i - 1)=3^ry+F_q(m)$, yields successive decompositions
	\[
	i = 2j \text{ or } 2j + 1,\quad
	j = 2k \text{ or } 2k + 1,\quad
	k = 2l \text{ or } 2l + 1,\ \ldots,\ 
	x = 2y \text{ or } 2y + 1.
	\]
	Consequently, we may write
	\[
	4i - 1 = 2^{s} y + \sum_{n_i \ge 2} 2^{n_i} + 3
	= 2^{s} y + m,
	\]
	which implies
	\[
	m \equiv 3 \pmod{4}.
	\]
	
	\medskip
	\noindent\textbf{Backward direction ($m \mapsto q$).}
	Let $m$ be any integer satisfying $m \equiv 3 \pmod{4}$. By
	Lemma~\ref{lem:3mod4}, $m$ admits a representation of the form
	\[
	m = \sum_{n_i \ge 2} 2^{n_i} + 3.
	\]
	Writing $4i - 1 = 2^ry + m$ and using the powers $2^{n_i}$ in the above decomposition,
	we recover successive relations
	\[
	i = 2j \text{ or } 2j + 1,\quad
	j = 2k \text{ or } 2k + 1,\quad
	k = 2l \text{ or } 2l + 1,\ \ldots,
	\]
	which uniquely determine an admissible sequence
	\[
	q = 11\cdots
	\]
	can be obtained via the inverse Collatz process and the mapping $F_q(4i - 1)=3^ry+F_q(m)$.
	
	\medskip
	Since the mappings $q \mapsto m$ and $m \mapsto q$ are inverses of each other,
	the correspondence between admissible sequences $q$ and integers
	$m \equiv 3 \pmod{4}$ is bijective.
\end{proof}


\begin{theorem}\label{thm:exhaustive}
	For any integer $i$, we have
	\[
	4i - 1 \equiv 3 \pmod{4}, \quad \text{and} \quad F_q(4i - 1) \neq 4i - 1.
	\]
\end{theorem}

\begin{proof}
	From above, we have
	\[
	4i - 1 = 2^{s} y + m,
	\]
	and
	\[
	F_q(4i - 1) = 3^{r} y + F_q(m),
	\]
	where $y$ is any integer.
	
	Suppose, for the sake of contradiction, that
	\[
	4i - 1 = F_q(4i - 1).
	\]
	Then it follows that
	\[
	m = F_q(m),
	\]
	and hence
	\[
	2^{s} y = 3^{r} y.
	\]
	Since $y \neq 0$, we obtain
	\[
	2^{s} = 3^{r}.
	\]
	
	This is impossible, because no power of $2$ can equal a power of $3$. Therefore,
	\[
	F_q(4i - 1) \neq 4i - 1,
	\]
	which completes the proof.
\end{proof}

\textbf{Remark:} This implies that for any finite sequence, the existence of a loop is impossible.



%\section{Stopping Number $F_q (m_c)$ and Cycle Number $m_c$}
\section{Bounds of Stopping Number and Cycle Number}


Suppose that there exists a cycle in the sequence, i.e. $$F_q (m_c )=m_c.$$ Using equation \eqref{formula}, we have $$\dfrac{3^rm_c+\sum_{i=1}^r3^{r-i} 2^{1(i)-1}}{2^s}=m_c.$$
Thus, 
\begin{equation*}
	m_c=\dfrac{\sum_{i=1}^r3^{r-i} 2^{1(i)-1}}{2^s-3^r}.
\end{equation*}

\begin{theorem}\label{thm:mc_low}
	The non-trivial cycle number $m_c$ is lower bounded by
	$$ \boxed {m_c\geq\dfrac{3^{r}-2^{r}}{2^s-3^r} }$$
\end{theorem}

\begin{proof}
Using \autoref{thm:min}, we immediately obtain

\begin{align*}
m_c=\dfrac{\sum_{i=1}^r3^{r-i} 2^{1(i)-1}}{2^s-3^r}\geq\min_{q} \dfrac{\sum_{i=1}^r3^{r-i} 2^{1(i)-1}}{2^s-3^r}
=\dfrac{3^{r}-2^{r}}{2^s-3^r}, 
\end{align*}
as claimed.

\end{proof}

\textbf{Remark:} Under the condition of Lemma~\ref{lem:length}, i.e., $ s = \left\lfloor 1 + r \cdot \log_2 3 \right\rfloor$, thus  $3^r<2^s$ holds, which makes the denominator $2^s-3^r$ positive.

For example, the value of expression
$$
\frac{3^r - 2^r}{2^s - 3^r}
\quad \text{for } r = 431{,}166{,}034{,}846{,}567 \text{ and } s = 683{,}381{,}996{,}816{,}440
$$
is approximately: \boxed{7.50 \times 10^{14}}. The values of $r$ and $s$ used here correspond to the maximum values listed in the reference \cite{El93}. Based on the results from \cite{Lu05}, the upper bound for the number of non-trivial cycles can be expressed as
 $$
\boxed { m_c < \dfrac{s \cdot 3^{r}}{2^s - 3^r}. }
 $$
 Substituting $r = 431{,}166{,}034{,}846{,}567$ and $s = 683{,}381{,}996{,}816{,}440$, this upper bound evaluates to approximately  \boxed{5.12 \times 10^{29}}.

Rewrite the stopping number as,
	$$F_q (m)=\dfrac{3^r m+\sum_{i=1}^r3^{r-i} 2^{1(i)-1}}{2^s}=\dfrac{3^r}{2^s}m +\sigma_q,$$
where $$\sigma_q=\dfrac{\sum_{i=1}^r3^{r-i} 2^{1(i)-1}}{2^s},$$
it is easy to see that the stopping number $F_q (m)$ consists of two terms, the first term includes starting number $m$ and is unrelated to the sequence $q$, while the second term doesn’t include starting number $m$, but it is only related to the sequence. The value of $F_q (m)$ is dominated by the first term. The second term determines the position of $F_q (m)$ between $m$ and $m/2$. Similar to the the cycle number, the lower bound of the second term $\sigma_q$ is determined by
$$\boxed{\sigma_q\geq\dfrac{3^{r}-2^{r}}{2^s}, }$$
while the upper bound based on the results of \cite{Lu05} is
$$\boxed{\sigma_q<\dfrac{s3^{r}}{2^s}, }$$
and from the observation data that are showing in \autoref{fig3}. Here, the starting numbers $m$ for \autoref{fig3}A, B, and C are from the positive integer of $3\pmod{4}$ start from 3 to 16,000,000 whose stopping time $s\geq4$. We can see that 
$$\boxed{\sigma_q<\alpha\dfrac{3^{r}-2^{r}}{2^s},}$$
where $\alpha$ is a constant dependent on the maximum value of $r$ in all sequences.

\begin{figure}[tb]
	\centering
	\includegraphics[scale=0.65]{fig3.png}
	\caption{Collatz process results for integers $m \equiv 3\pmod{4}$, ranging from 3 to 16,000,000. A. Second term $\sigma_q$ (blue dots) and its bounds in green and red versus $3^r/2^s$; B. Ratio $F_q (m)/m$ (blue dots) versus $3^r/2^s$; C. Starting number $m$ (blue dots), cycle number $m_c$ in green and its bounds in red and black versus $3^r/2^s$; D. Estimated bounds for the cycle number $m_c$ versus $3^r/2^s$ using extended $r$.}
	\label{fig3}
\end{figure}

From the computing results, we observed that \boxed{\alpha=max(r)/4}. Here $\alpha$=10 for  \autoref{fig3}A,B and C; $\alpha$=160 for \autoref{fig3}D with max of $r=640$. We also note that the estimated bond by Luca (2005) for cycle numbers in \cite{Lu05} has the same pattern as our calculated max of cycle numbers that plotted in the \autoref{fig3}D. The relationship between them can be regressed as  \boxed{y = 6.58 x -2.37}, which is plotted in the \autoref{reg_mc}.


\begin{figure}[tb]
	\centering
	\includegraphics[scale=0.55]{reg_mc.png}
	\caption{Regression of Estimated bound for the cycle number $m_c$ versus Calculated max values with $max(r)=640$.}
	\label{reg_mc}
\end{figure}


\autoref{fig3}A shows the second term of $F_q (m)$ (i.e. $\sigma_q$) and its bounds versus the ratio of $3^r$ and $2^s$. The \autoref{fig3}B shows that the non-trivial cycle number would be where $\dfrac{3^{r}}{2^s}$ is close to 1 if it existed. \autoref{fig3}B shows the ratio $F_q (m)/m$ versus ratio of $3^r$ and $2^s$. \autoref{fig3}C shows the starting number $m$ (blue dots) and the cycle number $m_c$ (green line) and its bounds (red lines). \autoref{fig3}D is showing the bounds of the cycle number with more values of $3^r/2^s$. The data in this figure was generated for $r=1$ to 640 (645 is the maximum number that R software can go with), then we apply $ s = \left\lfloor 1 + r \cdot \log_2 3 \right\rfloor$. The value of $s$ corresponding to $r$ is unique because of this restriction. Then we use the pair of $(r,s)$ to calculate the cycle number bound. We found that the minimum cycle number $m_c$ is greater than 978 when $(r,s)=(306,485)$ and $3^r/2^s\approx0.9989783$. It is not hard to see that the ratio $F_q (m)/m$ tends to be 1 as ratio $3^r/2^s$ goes to 1. Therefore, if there exists a cycle, the inequality $$\dfrac{3^r}{2^s} + \alpha\dfrac{3^{r}-2^{r}}{2^sm}\geq1$$
must be hold as $3^r/2^s$ tends to be 1. Thus, the cycle number $m_c$ must be less than a certain number, i.e.
\begin{equation} \label{mc_up}
 \boxed{m_c\leq\alpha\dfrac{3^{r}-2^{r}}{2^s-3^r},}
\end{equation}
which is also can be validate by the computing results in \autoref{fig3}C.

Finally, we have the following theorem.

%\begin{proposition}\label{prop:example}
\begin{theorem}\label{thm:final}
		
	The inequalities
	\[
	1-\epsilon_1 = \dfrac{1}{\alpha} +(1-\dfrac{1}{\alpha})\dfrac{3^r}{2^s} < \frac{F_{q}(m)}{m}  <\alpha+(1-\alpha)\dfrac{3^r}{2^s} = 1+\epsilon_2
	\]
	hold, and $\lim_{r \to \infty} \tfrac{F_{q}(m)}{m} \to 1$.
\end{theorem}

\begin{remark}
	This shows that any nontrivial cycle, if it exists, must be extraordinarily large and finely tuned, consistent with all known computational and theoretical bounds.
\end{remark}


\begin{proof}
	Building on the above discussion, and applying \autoref{thm:min} along with the established lower bound \autoref{thm:mc_low} and upper bounds \autoref{mc_up} for the cycle number, we obtain the following results:			
\begin{align*}
	\dfrac{F_q(m)}{m}&=\dfrac{3^r}{2^s} +\dfrac{\sum_{i=1}^r3^{r-i} 2^{1(i)-1}}{2^s m}\\ 
	&\geq\dfrac{3^r}{2^s} + \dfrac{3^{r}-2^{r}}{2^sm}\\
	&>\dfrac{3^r}{2^s} + \dfrac{2^{s}-3^{r}}{2^s\alpha}\\
	&=\dfrac{1}{\alpha} +(1-\dfrac{1}{\alpha})\dfrac{3^r}{2^s}\\
	&=1-\epsilon_1,
\end{align*}
and
\begin{align*}
\dfrac{F_q(m)}{m}&<\dfrac{3^r}{2^s} + \alpha\dfrac{3^{r}-2^{r}}{2^sm}\\
&\leq\dfrac{3^r}{2^s} + \alpha\dfrac{2^{s}-3^{r}}{2^s}\\
&=\alpha+(1-\alpha)\dfrac{3^r}{2^s}\\
&=1+\epsilon_2,
\end{align*}
respectively, so that $$ 1-\epsilon_1 = \dfrac{1}{\alpha} +(1-\dfrac{1}{\alpha})\dfrac{3^r}{2^s} < \frac{F_{q}(m)}{m}  <\alpha+(1-\alpha)\dfrac{3^r}{2^s} = 1+\epsilon_2.$$

Given $\alpha=max(r)/4$, then we have 
$\lim_{r \to \infty} \alpha \to \infty$, and we know that 
$\lim_{r \to \infty} \dfrac{3^r}{2^s} \to 1$, therefore,
$\lim_{r \to \infty} \tfrac{F_{q}(m)}{m} \to 1$. 

\end{proof}


The value of $1-\dfrac{3^{r}}{2^{s}}$ for $r = 431{,}166{,}034{,}846{,}567 \text{ and } s = 683{,}381{,}996{,}816{,}440$, is approximately: \boxed{1.33 \times 10^{-15}}.	

The Corollary 1.8 in \cite{Ev11} gives that				
$$1-\dfrac{3^r}{2^s}\geq(e\cdot s)^{-C},$$
where $C=e\cdot2^{3.5}30^5\log3\approx821013301.$ Thus, 
$$ (e \times 683{,}381{,}996{,}816{,}440)^{-821{,}013{,}301} \approx \boxed{10^{-7.1 \times 10^9}}. $$

Although this is unimaginably small—essentially zero for any computational purpose, $1-\dfrac{3^r}{2^s} \neq 0$, as $3^r \neq 2^s$.

\normalsize	

\section{Conclusion}

The results obtained in this paper indicate that the stopping time in the Collatz
process admits no finite upper bound. For admissible sequences $q$ with
parameters $(r,s)$, the ratio $3^{r}/2^{s}$ tends to $1$ as the sequence length
increases, yet it never equals $1$. Consequently, the quotient $F_{q}(m)/m$
approaches but never attains unity. This structural obstruction rules out the
possibility of a finite nontrivial cycle, since any such cycle would require an
exact equality between powers of $2$ and $3$. Combined with the lower and upper
bounds established in this work, these observations reinforce the conclusion that
all nontrivial cycles are excluded, and that the trivial cycle at $1$ is the only
admissible periodic behavior.

%    Bibliographies can be prepared with BibTeX using amsplain,
%    amsalpha, or (for "historical" overviews) natbib style.
% \bibliographystyle{amsplain}
%    Insert the bibliography data here.

\bibliographystyle{amsplain}
\begin{thebibliography}{10}
	
	%\bibitem {W} Wwikipedia, \textit{Collatz conjecture}, https://en.wikipedia.org/wiki/Collatz_conjecture.
	\bibitem {La11} J.C. Lagarias, \textit{The 3x+1 problem: An Annotated Bibliography (1963–1999) (sorted	by author)}, preprint, 2011. Available at \url{https://arxiv.org/abs/math/0309224}.
	\bibitem {La12} J.C. Lagarias, \textit{The 3x+1 Problem: An Annotated Bibliography, II (2000-2009)}, preprint, 2012. Available at \url{https://arxiv.org/abs/math/0608208}.
	%\bibitem {S} Stackexchange, \textit{3n+1 problem: fx(n)<n finitely many steps}, https://math.stackexchange.com/questions/3383810/3n1-problem-fxn-n-finitely-many-steps# 
	\bibitem {Si10} J. Simons and B. de Weger, Theoretical and computational bounds for m-cycles
	of the 3n + 1 problem, preprint, 2010. Available at \url{https://deweger.net/papers/
	[35a]SidW-3n+1-v1.44[2010].pdf}
	\bibitem {Ga81} L.E. Garner, \textit{On the Collatz 3n + 1 Algorithm}, Proceedings of The American Mathematical Society,  \textbf{82}(1), May 1981, 19-22.
	\bibitem {Mi01} T. Mimuro, \textit{On Certain Simple Cycles of the Collatz Conjecture}, SUT J. Math., \textbf{37}(2), 2001, pp. 79–89.
	\bibitem {Win17} M. Winkler, The algorithmic structure of the finite stopping time behavior of the 3x + 1 function, arXiv:1709.03385 [math.GM], 2017.
	\bibitem {El93} S. Eliahou, \textit{The 3$x$+1 problem: new lower bounds on nontrivial cycle lengths}. Discrete Mathematics, \textbf{118}(1-3), 1993, 45-56.
	\bibitem {Lu05} F. Luca, \textit{On the non-trivial cycles in Collatz’s problem}, SUT J. Math., \textbf{41}(1), 2005, 31–41.
	\bibitem {He23} C. Hercher, \textit{There are no Collatz m-cycles with $m<=91$}. Journal of Integer Sequences. \textbf{26}(3), 2023, Article 23.3.5.
	\bibitem{Ev11} J.H. Evertse, \textit{Linear forms in logarithms}, \url{https://www.math.leidenuniv.nl/~evertse/dio2011-linforms.pdf}, 2011.

	
	
\end{thebibliography}


%\section*{Statements and Declarations}
%
%	The author declare that no funds, grants, or other support were received during the preparation of this manuscript. The author has no relevant financial or non-financial interests to disclose. The author contributed to the study conception and design. Material preparation, data collection and analysis were performed by Daohang Sha. The datasets generated during and/or analysed during the current study are available from the corresponding author on reasonable request.


\section*{Supplemental Materials: R Programs}

Below is the R code used for generating all numbers $s, m, F_q(m),$ sequences $q$ in \autoref{fig:seqtree}, and many others based on $r$.

%\tiny
%\scriptsize
%\footnotesize
\small


\begin{lstlisting}
# function of even integer 
is.even <- function(x) {
	if (as.bigz(x) %% 2 == 0) TRUE
	else FALSE
}

# function to insert a substring in a string at certain position
insert_substr <- function(orig_str, substr, pos) {
   # Split the original string into two parts
   part1 <- substr(orig_str, 1, pos - 1)
   part2 <- substr(orig_str, pos, nchar(orig_str))

   # Concatenate the parts with the substring in between
   result <- paste0(part1, substr, part2)
   return(result)
}

# initialize                     
e<-NULL
n<- NULL
E<- NULL
r=2
s=floor(1+r*log23)
q <- "1100"
k=1
e[k]=2  # # of zeros of the sequence tail 
E[[r]]$e<-e
E[[r]]$q<-q

# not necessary
E[[r]]$m<-3
E[[r]]$d<-2
E[[r]]$r<-2
E[[r]]$s<-4

E[[1]]$r<-1
E[[1]]$s<-2
E[[1]]$e<-1
E[[1]]$q<-'10'
E[[1]]$m<-1
E[[1]]$d<-1

mk<-NULL
Fmk<-NULL

# start from 3:29 took 3.698423559202088206632 hours
for (r in 3:12){
   s=floor(1+r*log23)
   s1=floor(1+(r-1)*log23)
   d=s-s1
   ins <- ifelse(d == 2, "10", "1")

   pre_q <- E[[r-1]]$q
   e1<-E[[r-1]]$e
	
   n <- sum(e) # No. of residue classes in m(mod2^s)
	
   k=1
   # for each previous q's'
   for (u in 1:length(e1)){
      j=d
      q0<-pre_q[u]
      # get new q's  
      for (i in 1:e1[u]) {
         position <- s1-i+1
         q[k] <- insert_substr(q0, ins, position)  
         # cat("new q",k,"=",q[k], "\n")
         e[k] <- j
         j=j+1

         # get m and Fm from q created above
         # initialize the values and integer type (big z) 

         m=-1
         Fm=-1
         p=2
         t=2
			
         # use when s>50, about 3X slow
         # m=-big_1
         # Fm=-big_1
         # p=big_2
         # t=big_2
			
         for (i in 3:s) {
            step <- substr(q[k], i, i)
            if (step==0) {
               p=p #power of 3, max=r
               if (is.even(Fm)) {
                  Fm=Fm/2
                  m=m
               } else {
                  Fm=(3^p+Fm)/2
                  m=2^t+m
               }
            } else if (step==1) {
               p=p+1
               if (is.even(Fm)) {
                  Fm=(3*(3^(p-1)+Fm)+1)/2
                  m=2^t+m
               } else {
                  Fm=(3*Fm+1)/2
                  m=m
               }
            }
            t=t+1 # power of 2 max=s   
         }
         mk[k]<-m
         Fmk[k]<-Fm
			
         # Using big number package  
         # mk[k]<-as.numeric(m)
         # Fmk[k]<-as.numeric(Fm)
         k=k+1
       }
   }
   E[[r]]<-list(r=r,s=s,d=d,e=e,n=n,q=q, m=mk,Fm=Fmk)
}

\end{lstlisting}


\normalsize	
\end{document}



%-----------------------------------------------------------------------
% End of mcom-l-template.tex
%-----------------------------------------------------------------------
